\section{Configuration, executable identity and admission}
\label{architecture-admission}

The simulator accepts a document, but it executes a compiled plan. The YAML
reader first checks syntax, field names, types and ranges. Typed configuration
then identifies the selected grid, schemes, numerical parameters, output
policy and backend. The configuration compiler resolves the capabilities and
couplings those choices require. Finally, ordinary runtime admission decides
whether this particular combination can execute. These are separate questions:
a number can be syntactically valid yet lie outside an implemented physical
domain, and two individually available schemes can require an unfinished
coupling when selected together.

The authoritative path is visible in the
\href{../src/init/config/v1_loader.cpp}{typed loader},
\href{../src/core/runtime/configuration_compiler.hpp}{configuration compiler},
\href{../src/core/runtime/readiness.cpp}{readiness registry} and
\href{../src/core/runtime/nimbus_sim.cpp}{executable entry point}.
The registry distinguishes Engineering incomplete from Ready. The runtime
word \texttt{qualified} selects an admission policy; it does not assert that
a long campaign or external scientific review has occurred. Public Milbrandt
triple-moment selection remains guarded while its general source and live
thermal coupling are incomplete.

Studio does not maintain a second independent scientific model of which
configurations should work. Its frontend uses the executable's scheme catalog,
and its Rust preflight layer checks the source document and executable identity
again before launch. A browser snapshot provides the same catalog structure
when no native request is being made, but that snapshot is a generated input
with an authority fingerprint. An old snapshot must not silently authorize a
new executable. Conversely, a valid old run does not acquire new capabilities
merely because the currently installed Studio knows additional schemes.

Some narrow admission scopes bind exact canonical configuration digests. This
is stricter than matching a family name such as Metal or ZFP. The strict
serializer fixes the field inventory, defaults and representation; a membership
test then checks that canonical document against the declared scope. Separate
manifests bind raw source identity, serialized output-block identity and the
combined admission scope. Those hashes answer different questions. Formatting
the serializer can change its source hash without changing its serialized
block; changing the source-bound manifest then changes the hashes of downstream
combined manifests. Updating a block digest to stand in for a manifest digest
would be an identity error even if both strings are valid SHA-256 values.

The C++ catalog, TypeScript validator and Rust preflight therefore have
coordinated exact-scope checks. See
\href{../src/core/runtime/scheme_catalog.cpp}{catalog generation},
\href{../ui/src/lib/config/schemeCatalog.ts}{frontend validation} and
\href{../ui/src-tauri/src/commands/config_preflight.rs}{native preflight}.
These checks are not a substitute for numerical tests; they keep the tested
scope, selectable scope and user-visible scope from drifting apart. A metadata
refresh must preserve every non-hash field unless a separately justified scope
change is intended.

\section{Field storage, layout and ownership}
\label{architecture-storage}

\subsection{Contiguous fields and the cost of a traversal}

\href{../include/core/field/field3d.hpp}{Field3D} owns a contiguous array of
binary32 values. Although compatibility accessor names retain historical
radial/angular terminology, the current atmospheric configuration is Cartesian.
For dimensions $(n_x,n_y,n_z)$ the scalar address is
\[
  \operatorname{offset}(i,j,k)=(i n_y+j)n_z+k.
\]
The vertical index is contiguous. A vertical column sweep therefore visits
neighboring storage, while a horizontal stencil steps by a column or a slab.
This layout favors several column processes and permits simple full-volume
loops. It also means that a horizontal and a vertical stencil have different
cache behavior even on an isotropic physical grid.

State is organized as separate fields rather than an array of large cell
objects. A pressure operator can read density, theta and the required winds
without loading every microphysical number, radiation tendency or diagnostic
member attached to a cell. The cost is that an operator needing many fields
streams several arrays. For $N$ stored cells and $m$ complete Float32 fields,
the payload alone is $4mN$ bytes, before padding, geometry, snapshots, private
state, device copies or output. At $128\times128\times180$ cells, one field is
11.25 MiB. A seemingly minor extra read-and-write pass over ten such fields
represents approximately 225 MiB of logical traffic, even if caches or hardware
write policies change the actual memory-bus traffic.

The allocator aligns the base to 16 KiB and rounds the allocation size to that
alignment. This supports relevant host-memory import requirements and also
provides a well-aligned base for CPU access. It does not make every row aligned
to every SIMD width: the row stride still depends on $n_z$. Nor does alignment
prove that a Vulkan device accepts importing that allocation. Device memory
requirements, supported external-memory features, lifetime and synchronization
must still be checked. A successful aligned host allocation is only the first
part of such a contract.

Field dimensions are checked before allocation; copying a Field3D copies its
payload, while moving transfers its ownership. Raw pointer views do not extend
that lifetime. Aliasing is consequently an explicit operator concern. Some
helpers permit exact in-place input/output aliases, whereas shifted overlap
would overwrite unread input. A \texttt{restrict}-qualified vector loop is
only valid when its caller satisfies the non-aliasing promise. SIMD speed does
not justify relying on accidental overlap behavior. The
\href{../tests/core/test_simd_aliasing.cpp}{SIMD alias tests} exercise these
distinctions.

\subsection{Temporary capacity is reusable; contents are not automatically valid}

The \href{../include/core/field/field_pool.hpp}{field pool} groups reusable
buffers by shape. A normal acquisition clears recycled contents. An
overwrite acquisition skips that fill only because its caller promises to
write the complete required contents before reading them. Newly allocated
vector storage is still initialized by the container, so this optimization
removes repeated fills, not every allocation-time write. A poison mode fills
overwrite buffers with NaNs to expose incomplete initialization in tests.

The pool protects its bookkeeping with a mutex, but performs full-field fills
and copies after releasing that lock. This keeps independent shape requests
from serializing behind a large memory traversal. RAII guards return buffers
on normal exit or exception unwinding. A cached buffer is still allocated
memory: releasing a guard lowers live ownership, but does not necessarily
lower process RSS. Peak memory, current live bytes and retained reusable
capacity are different measurements.

This pattern matters especially on unified-memory systems. CPU scratch,
resident GPU resources, output snapshots and browser textures can compete for
the same physical capacity even though each subsystem reports a reasonable
local count. On a discrete system, the same workload can exhaust host RAM and
VRAM independently. A memory estimate must count the simultaneous ownership
graph, rather than multiply the number of prognostic variables by cell count
and assume everything else is negligible.

\subsection{Simulation state and nested contexts}

\href{../include/core/runtime/simulation_state.hpp}{SimulationStorageState}
owns the primary fields and associated runtime state. Views expose related
prognostic fields, base-state columns, tendencies, geometry and services to
operators. They carry access, not ownership. Compatibility entry points still
bind a global simulation instance; the design is not a claim that all legacy
code is free of globals. New code should use an explicit state view and the
correct level's runtime instead of recovering a backend or geometry through an
unrelated global.

A nested child has its own dimensions, physical domain, source history and
private scheme state. It cannot borrow a parent's buffer merely because the
shape matches. A moving grid also changes the physical interpretation of
array locations. Generation changes must invalidate stale geometry-dependent
or host/device reuse claims. This is one reason a domain move is a coordinated
transaction rather than an array shift applied independently to whichever
fields happen to be plotted.

\section{A timestep is a sequence of owned state transitions}

The \href{../src/core/orchestration/dynamics/dynamics.cpp}{dynamics owner}
connects slow tendencies, acoustic pressure/density and momentum updates,
boundary closure, retained transport fluxes, physical sources and the final
EOS projection. The mathematics guide derives the equations and their
locations. The architectural requirement is that every consumer knows which
stage its inputs represent. An old density, a newly projected pressure and
newly created cloud water are not one coherent stage merely because their
arrays have compatible shapes.

A typical bounded operator has a preparation phase, a candidate result and
a publication point. Preparation validates geometry, units, dimensions,
aliasing, domains and resources. Candidate work can populate scratch fields
or backend buffers. Publication changes the caller's authoritative fields
and the associated reports only after the result meets the owner's contract.
Failure atomicity is scoped to that transaction: it means that a rejected
operator does not leave a half-updated set of owned outputs. It does not
promise arbitrary rollback of the entire run or reverse changes made by an
unrelated earlier timestep.

Reports are outputs too. An operation changing momentum can also change
kinetic energy, boundary exchange and a diagnostic ledger. Updating the wind
while retaining the old report, or publishing a report before checking the
wind, creates inconsistent state. Similarly, microphysics primary water and
private number moments have a shared acceptance boundary. Checkpoint restart
is the separate mechanism for returning to an earlier completed model state.

Conservation diagnostics increase this coupling. A cell-centered internal
energy inventory and a C-grid face kinetic inventory need different geometry
and material inputs. The pressure-material reuse owner retains immutable
cell information while pressure changes; the velocity snapshot helper can
complete a snapshot without needlessly reconstructing unchanged thermal
parts. These are data-reuse changes with explicit preconditions, not changes
to physical equations. The relevant maintained sources include
\href{../src/core/orchestration/dynamics/pressure_cell_energy.hpp}{pressure
cell energy} and
\href{../tests/diagnostics/test_energy_budget.cpp}{energy-budget tests}.

The native fused acoustic path and the ordinary strict CPU path need not
traverse the same implementation. Profiling a similarly named helper is not
proof that the selected runtime calls it. An optimization must first identify
the actual admission and call path, then preserve the observable stage
sequence on that path. This avoids optimizing a dormant implementation while
claiming improvement in the active solver.

\section{ARM, x86 and x64 are several portability questions}
\label{architecture-cpu-portability}

\subsection{Instruction set is distinct from operating system and memory topology}

In this guide x64 means x86-64, the 64-bit x86 instruction-set family; ARM64
means AArch64. Neither term specifies the number of cores, vector width,
memory channels, cache layout, GPU, operating system or graphics driver.
An x86-64 host can use an integrated GPU with shared memory. An ARM64 host can
have a discrete accelerator. Selecting a memory policy from the CPU brand or
pointer width would conflate independent properties.

A native executable must match the target instruction set and operating-system
ABI. Object formats, calling conventions, system APIs and dynamic-library
dependencies also matter. A generic x86-64 Linux build is not a Windows binary
and is not an ARM64 binary. Replacing host-specific instructions does not
automatically make a new library ABI compatible with an older operating
system. A portable checkpoint similarly means a defined data format, not
permission to load it into an incompatible scientific configuration.

Nimbus's \href{../mk/options.mk}{build options} use
\texttt{-march=native -mtune=native} for ordinary source builds. These permit
the compiler to exploit the build host. \texttt{PORTABLE=1} removes those
host-native flags and records \texttt{NIMBUS\_PORTABLE\_BUILD};
\texttt{DISTRIBUTION=1} requires that portable choice and rejects conflicting
host-targeting flags. The resulting baseline remains the chosen compiler and
architecture's target. It is not a universal fat binary containing every ISA
variant. GCC explicitly distinguishes instruction-set selection from tuning
in its \href{https://gcc.gnu.org/onlinedocs/gcc/x86-Options.html}{x86 option
reference}.

On macOS, the build records a deployment target and distribution builds turn
unguarded availability into errors. Windows and Linux require their own
compiler/runtime dependencies and packaging checks. CPU-only, OpenMP-enabled
and serial builds also differ in capability. A serial build rejects an
explicit worker request greater than one instead of pretending that it
created a parallel team. The maintained
\href{SUPPORTED_PLATFORMS.md}{platform guide} states supported build evidence;
source branches for a platform are not evidence that every combination has
been exercised there.

\subsection{The actual SIMD implementation is narrower than the CPU's feature list}

\href{../include/util/simd_utils.hpp}{SIMD declarations} and
\href{../src/core/infra/simd_utils.cpp}{vector helpers} provide scalar, SSE,
AVX and Neon implementations for selected array operations. The implemented
paths process four Float32 values for SSE/Neon and eight for AVX. An AVX-512
compiler target deliberately uses the tested eight-lane AVX path in this
utility layer; there is no separate sixteen-lane implementation there.
Reported active width must describe that choice rather than advertising the
maximum feature named by the processor.

The code distinguishes compiled features from runtime host features. GCC/Clang
x86 capability queries use compiler builtins; AArch64 exposes the Neon path.
That reporting helps detect what is available, but the array implementations
are selected by compile-time branches. A runtime feature query cannot make a
whole translation unit compiled with unsupported instructions safe, and it
cannot create an AVX implementation absent from a generic binary. True
multiversioning would require separately compiled or target-attributed
implementations and a safe dispatch boundary. It must not be inferred from
the existence of a feature-report function.

The core solver also relies on compiler vectorization and OpenMP SIMD loops.
Therefore the helper width alone is not a census of all vector instructions in
a binary. A particular loop may stay scalar because of control flow, aliasing,
gather/scatter access, a function call or an unsupported transformation. The
useful evidence is the emitted code or vectorization report for the active
owner, accompanied by a correctness test. Wider vector instructions do not
guarantee proportional speedup when memory bandwidth or synchronization is
the limit.

Arm's \href{https://developer.arm.com/documentation/101028/0009/Advanced-SIMD--Neon--intrinsics}{Advanced SIMD reference}
describes fixed-width Neon operations. Nimbus's current helper does not become
an SVE or SME implementation merely by running on a newer ARM machine.
Likewise, generic x86-64 support and native AVX tuning are distinct tested
artifacts. Portable correctness should be checked with an actual baseline
object set, not by relinking a new entry point against old native objects.

\section{Threads, caches and deterministic budgets}

\subsection{Worker count is a policy input}

Nimbus uses OpenMP for shared-memory CPU execution. The CLI gives an explicit
\texttt{--threads} request priority and disables dynamic team sizing for that
request. Otherwise an existing \texttt{OMP\_NUM\_THREADS} remains authoritative.
Without either request, topology detection can prefer a detected performance
core class or physical cores on an SMT host, bounded by the runtime's available
thread count. It does not impose one hardcoded workstation thread count.
See \href{../src/core/infra/hardware_info.cpp}{hardware detection} and the
thread override in \href{../src/core/runtime/nimbus_sim.cpp}{the CLI}.

The detection code uses platform facilities: Apple performance-level counts,
Linux processor/package/core identifiers, and Windows processor topology.
These are observations, not a guarantee of exclusive access to those cores.
Container CPU sets, scheduler restrictions and an OpenMP runtime's processor
group behavior can constrain execution. A topology-derived default is a
starting policy rather than an automatic optimum for every mesh and scheme.

SMT adds logical execution contexts sharing core resources. Doubling logical
workers does not double floating-point units or memory channels. On a
bandwidth-heavy stencil, extra workers may increase competition; on a smaller
or latency-limited workload they may behave differently. Heterogeneous
performance and efficiency cores pose another scheduling question: one slow
participant can lengthen a barrier wait even if aggregate CPU utilization is
high. These effects are reasons to measure, not reasons to forbid a topology.

Nimbus records OpenMP binding and place information but does not install a
universal affinity policy. \texttt{OMP\_PLACES} and \texttt{OMP\_PROC\_BIND}
remain runtime controls. The OpenMP
\href{https://www.openmp.org/spec-html/5.1/openmpse62.html}{places} and
\href{https://www.openmp.org/spec-html/5.0/openmpse52.html}{binding}
definitions explain that these select where threads may run and how they
remain associated with places. A logged request is weaker evidence than the
runtime's actual placement; some environments cannot honor a requested policy.

\subsection{NUMA is a remaining measurement and placement concern}

The inspected implementation does not provide an explicit NUMA allocator,
\texttt{mbind}/\texttt{set\_mempolicy} placement policy or automatic
cross-socket work decomposition. Host vectors use ordinary allocation, and
their initialization and subsequent access interact with the operating
system's memory-placement policy. Consequently, a physical-core count alone
does not establish that a large array is close to every worker that reads it.
It would be inaccurate to describe Nimbus as automatically NUMA-optimized.

For a new multisocket or chiplet host, useful questions are which workers
touch which slabs, where pages are placed, which cache levels are shared, and
whether the working set saturates local or interconnect bandwidth. Changing
first-touch behavior or affinity is a separate candidate requiring matched
measurements. The current scalar layout and pool capacity reuse are useful
foundations, but they do not answer these topology questions by themselves.

\subsection{Reduction order is a numerical algorithm}

Floating-point addition is not associative. Dividing a sum among a different
number of workers usually changes its grouping, even if each cell contribution
is identical. Nimbus's deterministic diagnostic path instead forms fixed
horizontal slabs, visits each slab in a prescribed order, and merges slab
partials in increasing index with compensated accumulation. OpenMP can assign
the independent slabs to different workers without changing that partition.

For an accumulator $s$ and an added value $x$, compensation retains an estimate
of the low-order portion lost when $s+x$ is rounded. The implementation uses
a magnitude-dependent correction and returns the main sum plus the correction.
This improves the sum while fixed ordering supplies thread-count reproducibility.
Compensation alone would not fix an arbitrary parallel reduction tree.
The contract in
\href{../include/diagnostics/deterministic_reduction.hpp}{deterministic reductions}
is explicitly for one build target; it is not a blanket bitwise promise across
all libraries, architectures and source expressions.

Regular SIMD reduction helpers can use another grouping. They must not be
substituted for a deterministic scientific budget merely because their scalar
return type matches. Likewise, comparing two simulations requires identifying
whether equality means exact stored bytes, fixed-build repeatability, an
operator-specific tolerance, or a physical error bound. These are different
tests with different failure meanings.

\section{Backend dispatch, kernel ABI and memory admission}

The \href{../include/compute/backend_interface.hpp}{backend interface}
exposes operation families rather than an untyped function that accepts any
buffer list. Family availability, typed dispatch fields and validation define
what can execute. An absent family is unavailable; the owning runtime decides
whether that is a permitted fallback or a native-required failure. The backend
receives its runtime configuration before initialization so that independent
simulation contexts need not recover device selection through globals.

Each shader has a stable kernel identity and a manifest contract covering
resource bindings, constants, precision and availability. Host constants use
explicit integer widths, Float32 fields, padding and size assertions in
\href{../include/compute/kernel_abi.hpp}{the kernel ABI}. A C++ compiler's
ordinary object layout is not automatically a shader language's uniform or
push-constant layout. The explicit structures and manifest checks prevent an
apparently successful dispatch from interpreting an integer dimension as the
bits of a different parameter.

Ordinary native builds consume the committed SPIR-V artifacts after checking
the exact shader-source and included-GLSL hashes, artifact identities,
division-workaround variants, generated variant header and compiled kernel
manifest. They do not silently regenerate numerical kernels with whichever
shader optimizer happens to be installed on the build host. This matters
across machines: optimizer versions can remove unused push-constant padding,
and older SPIR-V tools may not understand an extension present in a retained
kernel. The \href{../scripts/dev/shaders/check_artifacts.py}{asset check}
rejects stale inputs; the separate
\href{../scripts/dev/shaders/verify_spirv.py}{rebuild/validation workflow}
requires a compatible toolchain and reviewed artifacts. Identity verification
preserves the selected implementation; it does not replace device admission,
dispatch parity or scientific validation.

The \href{../include/compute/dispatch_graph.hpp}{dispatch graph} records
read-after-write, write-after-read and write-after-write dependencies between
named field uses. It is a hazard description, not an executable resource
binder. A structurally valid graph still needs the correct buffers, extents,
constant block, kernel and backend synchronization. This distinction keeps a
graph validation pass from being misreported as successful GPU execution.

Memory admission is also backend-resolved. The largest individual field must
fit the device's storage-buffer limit, while simultaneous resident fields
must fit an available heap budget. Those are independent constraints. The
current dynamics envelope uses 51 field-equivalents per simultaneously resident
level, including retained names and transients; it is a conservative planning
envelope, not the exact number of prognostic variables. Specialized wider
transactions retain their own limits. Saturating arithmetic reports overflow
as invalid admission rather than wrapping a huge request into a small one.
See \href{../include/compute/compute_memory_admission.hpp}{memory admission}
and \href{../tests/compute/test_memory_admission.cpp}{its tests}.

If $B_f=4n_xn_yn_z$ is one field payload, the simple envelope is approximately
$51B_f$ per level before interpreting host/device duplication. A hierarchy
must sum all levels that can be resident together. A device-local heap size is
not necessarily its currently free budget; Nimbus distinguishes a known heap
from a live budget when the backend exposes one. Admission can be conservative,
but it must not invent VRAM capacity from a CPU's system-memory report.

\section{Vulkan transfers, execution and failure containment}
\label{architecture-vulkan}

\subsection{Choose memory from reported types, not a GPU name}

\href{../vulkan/include/compute/vulkan_memory_plan.hpp}{The Vulkan memory plan}
inspects compatible memory types and their heaps. It identifies a primary
device-local heap, a working allocation type and a host-visible staging type.
For staging it prefers cached host memory, which matters for CPU readback.
Host coherence and host caching are different flags: coherence concerns
visibility operations, while caching affects CPU access behavior.

A type that is both host-visible and device-local can describe genuinely
shared memory, or a CPU-visible aperture into discrete VRAM. Resizable BAR
does not turn the latter into an integrated-memory machine. Nimbus keeps a
discrete device in an aperture/staged classification even when the visible
aperture spans the full device heap. Direct shared mode requires compatible
reported topology; \texttt{auto}, \texttt{direct} and \texttt{staged} controls
select among supported plans rather than bypassing their requirements.
The underlying distinction is described in the
\href{https://docs.vulkan.org/guide/latest/memory_allocation.html}{Khronos
memory-allocation guide}.

The buffer pool retains capacity for repeated dispatches. In direct shared
mode a device buffer is mapped and the staging slot can be a placeholder.
In staged mode the pair contains host-accessible staging and device-local
storage. Batch allocation can share a backing allocation while retaining
separate buffer handles. Descriptor updates still bind each logical resource
with its correct range. Retaining capacity reduces allocator work; it does
not remove the obligation to upload newly changed data.

\subsection{The staged sequence has two kinds of copies}

A representative synchronous operation follows this sequence. It is a
description of ownership stages, not pseudocode authorizing arbitrary
barrier removal:

\begin{verbatim}
validate dispatch and obtain owned candidate buffers
copy host input arrays into mapped staging
flush noncoherent host writes when required
record staging -> device copies
barrier: transfer writes become shader-readable
record compute dispatches and inter-kernel dependencies
barrier: shader writes become transfer-readable
record device -> staging copies
submit command buffer; wait for confirmed completion
invalidate noncoherent host mappings when required
validate result; copy/publish accepted host outputs
\end{verbatim}

The CPU copy into staging and the GPU transfer into VRAM are separate work.
Similarly, downloading to staging and copying into the solver's authoritative
host arrays are separate work. Telemetry distinguishes host upload/download
bytes from device upload/download bytes. Adding those quantities without
stating the boundary can double-count one logical field movement, while
omitting either can hide the actual cost paid by the process.

A pipeline barrier describes ordering and memory visibility between producer
and consumer accesses. A fence lets the host learn that submitted work has
completed. A mapped-memory flush makes noncoherent host writes available; an
invalidate makes completed device writes visible to a noncoherent host mapping.
None is a synonym for the other. Host-coherent memory removes particular cache
maintenance calls, not the need to wait before reading a result or overwriting
in-flight input. Khronos's
\href{https://docs.vulkan.org/guide/latest/synchronization_examples.html}{synchronization examples}
illustrate these distinct dependencies.

Nimbus aligns noncoherent ranges to the device's atom-size rules and avoids
processing one shared allocation repeatedly through multiple handles.
The authoritative implementation is
\href{../vulkan/src/compute/compute_backend_vulkan.cpp}{the Vulkan backend}.
A change that removes a copy must also update the relevant authority and
visibility transitions; deleting \texttt{memcpy} from one branch is not a
complete residency optimization.

\begin{figure}[htbp]
\centering
\begin{tikzpicture}[
  box/.style={draw=NimbusBlue,fill=NimbusLight,rounded corners=2pt,
              text width=3.2cm,minimum height=12mm,align=center,font=\small},
  arrow/.style={-{Latex[length=2mm]},draw=NimbusBlue,thick},
  node distance=18mm and 9mm]
  \node[box] (host) {Authoritative host fields\\version $v$};
  \node[box,right=of host] (staging) {Mapped staging\\CPU-visible transfer bytes};
  \node[box,right=of staging] (device) {Device candidate\\bound buffers and kernel};
  \node[box,below=of host] (commit) {Validate and publish\\host version $v+1$};
  \node[box,below=of staging] (readback) {Completed readback\\visibility established};
  \node[box,below=of device] (done) {GPU output\\fence-confirmed completion};
  \draw[arrow] (host)--node[above,font=\scriptsize]{copy}(staging);
  \draw[arrow] (staging)--node[above,font=\scriptsize]{copy}(device);
  \draw[arrow] (device)--node[right,font=\scriptsize]{dispatch}(done);
  \draw[arrow] (done)--node[below,font=\scriptsize]{copy}(readback);
  \draw[arrow] (readback)--node[below,font=\scriptsize]{copy}(commit);
\end{tikzpicture}
\caption{A representative staged-memory compute transaction. Direct shared
memory can remove the separate staging transfer, but still requires ordered
access and an authoritative version transition. The placement of validation
and commit is an operator contract; uncertain GPU execution cannot release
the resources it may still reference.}
\label{fig:architecture-staged-memory}
\end{figure}


\subsection{An uncertain submission cannot release its buffers}

The current compute owner records and submits one shared command buffer at a
time. A successful queue submission means the queue accepted work, not that
its outputs are ready. Completed-operation counters advance after fence
confirmation and required visibility handling. Kernel counts are associated
with that completed submission, so a failed or abandoned command is not
reported as successful native science.

If submission or waiting leaves execution uncertain, the backend is poisoned
and referenced resources remain pinned until shutdown. Reusing a pool slot
after a timeout could let the GPU write into a later unrelated operation.
The pure \href{../vulkan/include/compute/vulkan_submission_lifecycle.hpp}{submission
lifecycle} models this independently of real handles so fault tests can cover
the decision. A host fallback is possible only where the enclosing runtime
allows it and retains a separate valid host state; native-required admission
does not become optional because a device failed.

\section{Residency is a versioned authority graph}

A field allocation can exist without containing the current field. Nimbus's
\href{../include/compute/compute_residency.hpp}{residency graph} records a
logical version, host version and device version together with host identity
and extent. A host write advances the logical version and host version; a
device write advances the logical version and device version. Synchronization
copies a known version to the other side. Reuse is valid only when the side's
version equals the authoritative version and identity/extent still match.

For example, suppose both copies hold version 7. A host boundary update makes
the host version 8 while the device remains at 7. A subsequent kernel cannot
skip upload merely because its input pointer and allocation handle are
unchanged. Conversely, after a device-only update produces version 9, a host
diagnostic must not read version 8 just because the host pointer remains valid.
The state machine exposes these mismatches; it does not itself execute a fence
or inspect the pointed-to bytes.

\href{../include/compute/resident_field_store.hpp}{ResidentFieldStore}
combines that authority graph with named allocation claims. Replacing allocator
callbacks while claims are live is rejected. Batch acquisition publishes no
partial set if a candidate allocation fails, and the callbacks must outlive
the claims that use them. The store is not thread-safe: the backend must
serialize its transitions. Its destructor intentionally does not call an
allocator that may already have been destroyed; owners needing individual
release use the explicit clear path while the allocator is alive.

Generation identity is especially important for restart, grid movement,
resizing and nested contexts. Pointer reuse by a host allocator is not proof
that an old device value still represents the same simulation field. The
\href{../tests/compute/test_residency.cpp}{residency tests} cover state and
allocation behavior, while full runtime tests must additionally prove that
every actual writer announces its transition. A correct cache state machine
cannot compensate for a producer that silently mutates storage behind it.

\section{Metal follows the same ownership questions with different APIs}

Metal exposes a device, command queue, buffers, pipeline states and command
encoders rather than Vulkan's object model. Nimbus stores their session
identity in
\href{../metal/src/compute/internal/metal_device_context.hpp}{MetalDeviceContext}.
Its synchronous finish path commits the command buffer, waits for completion,
checks the final status and only then advances completion telemetry. GPU start
and end timestamps describe the completed command; host submit/wait time is
a different interval.

Shared-storage buffers make memory CPU- and GPU-accessible on the applicable
device, but a host vector and a Metal buffer created with
\texttt{newBufferWithBytes} are still different owned objects. Creating that
buffer copies bytes. Copying its contents into a solver field after completion
is also real work. Some operators retain resident named resources; others
still allocate or copy through their ordinary paths. A blanket statement that
Metal avoids all transfers would obscure the actual implementation.

Apple's \href{https://developer.apple.com/documentation/metal/choosing-a-resource-storage-mode-for-apple-gpus}{storage-mode guide}
distinguishes CPU-visible shared resources from GPU-private resources. A
private resource may require explicit transfer access when a CPU needs its
data, while shared access still needs correct execution ordering. A private
storage experiment does not become the production policy merely because a
local benchmark can create such resources. Nimbus's
\href{../metal/src/compute/backend/advection.inc}{advection owner} and
\href{../metal/src/compute/internal/metal_resident_state.hpp}{resident state}
show concrete copy and reuse paths.

The backend also checks kernel coverage against the canonical manifest. A
kernel demanding native Float64 arithmetic is reported as hardware-limited
in the current Metal coverage contract; it is not secretly replaced by a
Float32 kernel and called equivalent. Other scientific kernels use explicitly
maintained compensated Float32 formulations. Availability of one family says
nothing about an unrelated family's numerical requirements.

\section{Sharding a simulation across multiple GPUs}
\label{architecture-multigpu}

This section is a design analysis of work Nimbus would need. It does not
describe an implemented distributed solver. The current backend configuration
selects one device index for a compute runtime. The
\href{../vulkan/include/compute/vulkan_device_context.hpp}{Vulkan context}
owns one selected physical device, logical device and compute queue; the
\href{../metal/src/compute/internal/metal_device_context.hpp}{Metal context}
owns one device session and command queue. There is no maintained MPI domain
decomposition or complete multi-device timestep owner in the ordinary runtime.
Multiple runtimes, nested domains, CUDA-assisted kernels and remote SSH
execution do not, individually or together, establish that capability.

\subsection{What is being divided?}

There are three substantially different designs. Independent simulations can
run on different devices without exchanging atmospheric state. Spatial
decomposition instead divides one atmosphere into owned subdomains, which
must exchange the state needed by neighboring numerical operators. Assigning
different physics packages to different GPUs is a third design: for example,
dynamics on one device and radiation on another. That arrangement must move
inputs and return tendencies at their declared cadence. It may duplicate much
of the state and adds a coupling dependency, so it should not be assumed to
solve a memory-capacity problem.

For an initial spatial implementation, horizontal tiles with complete vertical
columns are a reasonable candidate. They keep column radiation, boundary-layer
solves, vertical sedimentation and a column-local implicit solve on one owner.
This is a proposed starting restriction, not proof of support for every current
column formulation. It can produce poor aspect ratios or leave too few tiles
for a large machine. Splitting vertically later would require explicit
distributed column algorithms or gather/scatter phases, including their
memory and synchronization costs. Terrain and moving nests make ownership
more complex even when columns remain whole.

\begin{figure}[htbp]
\centering
\begin{tikzpicture}[
  box/.style={draw=NimbusBlue,fill=NimbusLight,rounded corners=2pt,
    text width=4.9cm,minimum height=17mm,align=center,font=\small},
  arrow/.style={-{Latex[length=2mm]},draw=NimbusBlue,thick}]
  \node[box] (a) at (0,0) {GPU A: owned horizontal tile\\complete vertical columns};
  \node[box] (b) at (7.2,0) {GPU B: owned horizontal tile\\complete vertical columns};
  \draw[arrow] ([yshift=3mm]a.east)--node[above,font=\scriptsize]{stage $s$ halo}([yshift=3mm]b.west);
  \draw[arrow] ([yshift=-3mm]b.west)--node[below,font=\scriptsize]{stage $s$ halo}([yshift=-3mm]a.east);
  \node[box] (stage) at (3.6,-2.4) {Validate the shared stage\\global timestep and owned budgets};
  \draw[arrow] (a.south)--(stage.north west);
  \draw[arrow] (b.south)--(stage.north east);
  \node[box,below=8mm of stage] (commit) {One accepted endpoint\\consistent output and checkpoint};
  \draw[arrow] (stage)--(commit);
\end{tikzpicture}
\caption{Proposed two-tile solver decomposition. Halo exchange can use a
validated direct path or explicit host staging. Arrows describe data
dependencies, not a mandatory global barrier after every local operator.
Each tile owns its physical cells once; copied halos carry no additional
atmospheric inventory. This is a design diagram, not a current Nimbus feature.}
\label{fig:architecture-multigpu}
\end{figure}


\subsection{Own physical cells and faces before assigning buffers}

Each tile needs global physical bounds, local storage bounds, halo widths,
neighbors and a stable mapping from global to local indices. A halo is a
readable copy of another tile's state, not additional atmospheric volume.
Inventory reductions must count owned cells once. On a C-grid, a velocity or
mass-flux face shared by two tiles also needs a canonical owner or an equivalent
reconciliation rule. Independently evaluating slightly different fluxes on
both sides can create an interface source even when each local kernel passes
its own conservation test.

For one mass flux $\Phi_f$ oriented from tile $A$ to tile $B$, conservative
exchange requires
\[
 \Delta M_A=-\Delta t\,\Phi_f,\qquad
 \Delta M_B=+\Delta t\,\Phi_f.
\]
The same face, time interval, stage weight and transported carrier must be used
on both sides. Water mixing ratios, number per dry mass and number per volume
cannot share an unspecified scalar convention. Geometry, metric coefficients,
held forcing and scheme-private fields require exchange wherever their
consumers need them. A fixed list of visible plot fields is not a sufficient
halo protocol.

Messages should identify the grid or hierarchy generation, physical interval,
integrator stage, field and location, source and destination ownership, extent
and precision. A completed transfer with the wrong stage is stale data. Halo
width must follow the largest actual dependency footprint of the selected
operator, including reconstruction and corner dependencies; it cannot be
guessed from the nearest-neighbor appearance of the final divergence.

Physical boundary conditions apply at the atmosphere's external boundary.
An internal partition face must not become a wall, sponge or independent open
boundary. Periodic neighbors must close the global topology. A processor
partition on a fixed grid is also different from a coarse/fine nest: the
latter needs metric-aware restriction, prolongation and flux reconciliation,
not merely a copy between equal-resolution halos.

\subsection{Exchange at the stage that consumes the data}

A portable first implementation could establish this sequence for each
operation with neighbor dependencies: complete the producing work, pack the
required owned boundary state, start its transfers, compute independent
interior cells, complete and unpack the receives, then compute the boundary
region. The new endpoint becomes visible only after all required regions
and validations complete. Packing, transfer and unpacking each retain buffer
ownership until their respective completion events.

This sequence must be derived from the actual split integrator. Acoustic
substeps may require pressure and face-wind exchange several times during one
slow step. Runge--Kutta stages consume their own stage state. Exchanging only
after a complete outer timestep is generally insufficient. A global CFL
decision must use the most restrictive participating tile; a local reduction
does not establish the safety of the global timestep. Iterative diagnostics
or solvers can add global reductions beyond those required by local stencils.

The current \href{../include/compute/dispatch_graph.hpp}{dispatch graph}
describes field hazards, but is not a distributed execution engine. It would
need region and device ownership, explicit exchange edges and completion
dependencies. Likewise, the
\href{../include/compute/compute_residency.hpp}{residency contract} would need
to distinguish an owned interior version from a halo version for a particular
stage. Marking an entire array ``current on the device'' cannot express that
one boundary slab has not arrived.

\subsection{Portable communication first, optional direct paths second}

The reference transport should work with explicit device-to-host transfer,
owned host send/receive buffers and host-to-device transfer. This gives a
testable route without assuming peer access. On a single machine the exchange
can be managed within one process or across explicitly owned processes.
Across machines, MPI is a plausible transport and collective layer. Its
\href{https://www.mpi-forum.org/docs/mpi-4.1/mpi41-report/node235.htm}{nonblocking
neighborhood operations} provide an interface for neighbor exchange, but
posting a nonblocking operation is not proof of overlap or completion. Progress,
thread support and the application's wait schedule must be measured and
validated. The transport does not know Nimbus's field units or stage semantics.

Direct device transfers are optional optimizations of the same exchange
contract. Vulkan device groups are enumerated compatibility groups; they do
not promise that arbitrary installed GPUs can be treated as one memory pool.
The \href{https://docs.vulkan.org/guide/latest/extensions/device_groups.html}{Khronos
device-group guide} and
\href{https://docs.vulkan.org/refpages/latest/refpages/source/vkGetDeviceGroupPeerMemoryFeatures.html}{peer-memory
query} describe the relevant capability checks. Separate logical devices
instead require a supported sharing/transfer mechanism with explicit lifetime
and synchronization rules. A queue barrier or fence on one independent device
does not by itself establish that another device can consume the bytes.

Metal similarly requires using the devices actually exposed by the system.
Apple documents
\href{https://developer.apple.com/documentation/metal/multi-gpu-systems}{multi-GPU
systems} and
\href{https://developer.apple.com/documentation/metal/transferring-data-between-connected-gpus}{transfers
between connected GPUs}. Availability depends on hardware and API support;
multiple packages or dies inside a product do not imply multiple independently
selectable Metal devices. Unified CPU/GPU memory on one device does not establish
peer access to another computer or accelerator.

``GPU-aware MPI'' also needs a specific memory API, transport and build. For
example, \href{https://docs.open-mpi.org/en/v5.0.8/tuning-apps/networking/cuda.html}{Open
MPI's CUDA support} describes supported CUDA-buffer paths. It does not certify
arbitrary Vulkan allocations or Metal buffers. Nimbus would need a validated
interoperability path or host staging for those owners. An unavailable direct
path may select an explicitly supported staged transport; it must not silently
change the numerical operation or violate a native-execution requirement.

\subsection{Why more memory does not imply linear speedup}

Let an illustrative cubic owned tile have side $n$, halo depth $h$, $F$
communicated scalar-equivalents and $b$ bytes per scalar. Its owned payload is
$bFn^3$. A complete surrounding halo occupies
\[
 B_{\mathrm{halo}}=bF\big[(n+2h)^3-n^3\big]
 =bF(6hn^2+12h^2n+8h^3).
\]
If only the six face slabs are communicated in an exchange, their leading
payload is $6bFhn^2$; corners, staged directional exchanges and non-scalar
locations change the exact count. This is a teaching model for an interior
tile, not the traffic of a particular Nimbus kernel. As a fixed global problem
is divided into more pieces, the surface-to-volume ratio grows. Full-column
horizontal tiles have a different geometric count but the same concern.

For message startup cost $\alpha$, effective transfer bandwidth $\beta$,
$m$ messages and payload $B$, a simple model is
\[
 T_{\mathrm{exchange}}\simeq m\alpha+B/\beta
       +T_{\mathrm{pack}}+T_{\mathrm{unpack}}.
\]
With perfect interior/exchange overlap, and with boundary work and a collective
following those phases, a stage still takes roughly
\[
 T_{\mathrm{stage}}\gtrsim
 \max(T_{\mathrm{interior}},T_{\mathrm{exchange}})
 +T_{\mathrm{boundary}}+T_{\mathrm{collective}}.
\]
Real overlap can be worse: packing can contend for memory bandwidth, transfer
engines may share links, and small tiles can underfill a GPU. Frequent acoustic
exchanges can become latency-bound. A second board behind a narrow or shared
link may increase capacity while adding little useful throughput.

Every device must fit its own owned state, halos, stages, rollback and
communication buffers. The sum of nominal VRAM capacities is not an admission
test. Nor should every rank retain a full global host mirror: that would
preserve a memory bottleneck even after the device fields were partitioned.
NUMA placement should associate host staging buffers and communication threads
with the relevant device and network path. On heterogeneous GPUs, equal cell
counts need not give equal stage times. Cloud-dependent work and nested
placement can also make the slowest tile change during an experiment.
Dynamic balancing would require a state migration transaction and new
generations, not just moving an array pointer.

\subsection{Global diagnostics, restart and failure must agree}

Budget reductions should accumulate owned physical inventories with declared
precision and ordering. A reduction tree may change rounding when the number
of ranks changes. Fixed-decomposition repeatability, decomposition-independent
bitwise equality and bounded scientific agreement are distinct claims. They
need distinct tests; ordinary parallel summation does not promise all three.
Random perturbations should have a defined mapping to physical identities and
time if results are intended to be independent of rank placement.

A checkpoint needs a globally consistent accepted endpoint: all owned primary
and private fields, geometry, held-source clocks, histories, configuration and
publication identity. Repartitioned restart needs a mapping from saved physical
ownership to the new layout; it cannot depend on finding the same GPU index.
Partitioned output should publish a manifest only after all required pieces
are durable, and readers should reconstruct the logical field without treating
rank order as a coordinate system. Gathering the whole atmosphere on one rank
for every frame can become a new memory and I/O bottleneck. Existing NPY,
netCDF and ZFP support does not by itself supply this distributed publication
protocol, and this design does not authorize removal of NPY.

If one tile rejects a timestep, neighboring tiles cannot publish a mixture of
old and new endpoints. A first design should stop the distributed job on an
unrecoverable device or communication failure and resume from a validated
global checkpoint. Transparent continuation after losing a rank is additional
engineering, not a consequence of using MPI. Buffers involved in uncertain
device submissions must remain owned until completion is known or the
corresponding execution context has been safely retired.

\subsection{A sequence that can be implemented and reviewed}

Begin with the distributed state and exchange contract, including a one-tile
path that reproduces the existing admitted experiment. A two-tile CPU or
host-staged reference can then expose ownership errors independently of peer
transport. Manufactured fields should cross each partition orientation and
corner; interface inventories must cancel, and halos must never enter physical
budgets. Tests should mutate stage identities, interrupt transfers and reject
partial publication.

The next bounded target is a two-device, fixed-grid configuration with stated
physics and boundary conditions, selected through ordinary public configuration.
It needs whole-step one-device comparisons, native dispatch evidence, restart
on the same decomposition, rejected-input atomicity and memory accounting.
Changing the decomposition on restart, nonuniform hardware, terrain, private
microphysics state, held radiation, nesting and moving domains each extend the
contract and require their own completed coupling. They cannot be called
supported merely because the flat dry exchange works.

Only after correctness should direct peer transfers, overlap, larger device
counts and multi-node transports be assessed with strong- and weak-scaling
experiments. Strong scaling holds the physical numerical problem fixed; weak
scaling holds approximately the work per participant fixed and generally grows
the global problem. Report both completion time and resource cost, including
output and restart. Source code or a two-GPU speedup alone is not Ready, and
Ready remains separate from campaign validation and independent scientific
review. This staged design keeps the ordinary local workflow usable while
making distributed execution an explicit, testable extension.


\section{Floating-point semantics across CPU and GPU}
\label{architecture-numeric-portability}

\subsection{Storage precision, arithmetic precision and rounding points}

Binary32 has 24 bits of significand precision, including its implicit leading
bit for normal values; binary64 has 53. Their unit roundoffs near one are
approximately $2^{-24}$ and $2^{-53}$ under round-to-nearest. Storing the main
fields as Float32 therefore limits what a wider intermediate can publish.
If a source computes a tiny increment in double precision and then stores
\texttt{float(old + increment)}, that increment can disappear at the storage
boundary. It is not necessarily safe to move the cast earlier, because an
intermediate rounding can also change a threshold decision or cancellation.

Nimbus's dry EOS deliberately evaluates from stored Float32 density and theta
using wider arithmetic before storing pressure. Energy attribution distinguishes
ideal changes from water, theta and pressure storage effects. Backend work must
preserve the intended rounding points and constants. A tolerance patch that
makes one Linux failure disappear without explaining the changed operation
would conceal a numerical portability problem rather than solve it.

\subsection{FMA is useful precisely because it changes rounding}

A fused multiply-add evaluates $ab+c$ with one final rounding, whereas separate
operations compute a rounded product before the sum. Schematically,
\[
 \operatorname{fma}(a,b,c)=\operatorname{RN}(ab+c),\qquad
 \operatorname{RN}(\operatorname{RN}(ab)+c)
 \text{ need not be equal.}
\]
The distinction is especially visible when $ab$ and $c$ nearly cancel.
Compensated algorithms can intentionally use FMA to recover a product residual:
if $p=\operatorname{RN}(ab)$, then $e=\operatorname{fma}(a,b,-p)$ captures the
remaining portion under the algorithm's finite-range assumptions. Reassociating
that expression or replacing it with an ordinary multiply and subtract defeats
the construction.

Even a helper named \texttt{fma\_vectors} needs its concrete implementation
examined. Nimbus uses the fused AVX intrinsic when \texttt{\_\_FMA\_\_} is
compiled, separate multiply/add in the SSE and non-FMA AVX branches, and a
Neon fused intrinsic in its AArch64 vector branch. Scalar tails have their
own compiler contraction behavior. Thus a name alone does not establish
bitwise-identical evaluation on every architecture. Scientific kernels with
stronger requirements use their own explicit arithmetic contracts.

\subsection{Compensated pairs and device capability gates}

Several GPU kernels represent a more accurate intermediate as a high/low
Float32 pair. Addition and multiplication propagate low-order residuals with
carefully ordered fused operations. This increases useful precision for a
defined calculation, but it is not a universal replacement for IEEE binary64:
the exponent range, subnormal behavior, transcendental functions and final
storage conversion still need their own treatment. The
\href{../metal/include/compute/metal/shaders/terrain/mapped_scalar_math_msl.hpp}{mapped
scalar Metal helper} is generated from shared arithmetic sources rather than
being separately hand-tuned into a different equation.

Vulkan distinguishes ordinary Float32 kernels from those explicitly requiring
Float32 storage with Float64 arithmetic. The latter are gated on the required
Float64 rounding/preservation properties, strict Float64 FMA and integer
storage-conversion capability. Compensated families also require the relevant
Float32 signed-zero/NaN preservation and round-to-nearest controls. A device
advertising basic compute support or a high peak throughput number does not
thereby satisfy these scientific requirements.

Tiny values introduce another issue: devices can flush subnormal arithmetic
to zero. The maintained
\href{../vulkan/shaders/compute/terrain/compensated_storage.glsl}{storage helper}
scales tiny inventories before arithmetic and uses explicit integer-bit
rounding for the subnormal storage case. That is an algorithm-specific
representation technique with tests, not permission to add platform-specific
epsilon clamps to every failed comparison. Metal's retained compile policy
disables fast math for these owned shader sources; generator checks keep
translated arithmetic synchronized.

The practical validation sequence follows the contract: check invalid inputs
and output atomicity, compare the intended operator against its independent
reference, test native and generic CPU builds where applicable, then exercise
the real selected backend and memory policy. Exact saved-state comparison,
bounded numerical tolerance and scientific closure remain distinct evidence.
Neither a portable build nor a parity test alone proves that every tracked
configuration is Ready.
